\chapter{สรุปผล อภิปรายผล และข้อเสนอแนะ}

\section{สรุปผลการดำเนินงาน}

โครงงานนี้ได้พัฒนาแบบจำลองการเรียนรู้ของเครื่องสำหรับทำนายค่าฤทธิ์การยับยั้งเอนไซม์ SARS-CoV-2 Main Protease ในรูป \pICfifty\ แบบ In Silico โดยดำเนินงานครบทั้ง 9 ระยะตามแผน สรุปผลได้ดังนี้

\textbf{ด้านข้อมูล} รวบรวมข้อมูลจาก ChEMBL เป้าหมาย \texttt{CHEMBL4523582} จำนวน 4{,}867 รายการ กรอง assay เฉพาะที่วัด \CLpro\ เหลือ 2{,}889 รายการ (คัดออกร้อยละ 40.6 ซึ่งเป็น assay ของ nsp3, PLpro และ papain-like protease) ผ่านการทำความสะอาด 13 ขั้นตอน ได้ชุดข้อมูลมาตรฐาน \textbf{2{,}198 โมเลกุลที่ไม่ซ้ำกัน} ประกอบด้วยสาร active 1{,}453 ตัว intermediate 547 ตัว และ inactive 198 ตัว มีความหลากหลายเชิงโครงสร้างสูงโดยมี Murcko Scaffold ไม่ซ้ำ 976 แบบ

\textbf{ด้านแบบจำลอง} XGBoost ให้ประสิทธิภาพสูงสุดบนชุดทดสอบ โดยได้ \Rsq\ = 0.7034, RMSE = 0.5487, MAE = 0.3918 และ Spearman $\rho$ = 0.8415 ซึ่งเหนือกว่าแบบจำลองพื้นฐานอย่างชัดเจน (\Rsq\ = $-0.0000$) การทำนายร้อยละ 93.6 มีความคลาดเคลื่อนไม่เกิน 1 log unit เมื่อประเมินด้วยการแบ่งข้อมูลตามโครงสร้างแกนกลาง ค่า \Rsq\ ลดลงเหลือ 0.5127

\textbf{ด้านการวิเคราะห์ SAR} ระบุคุณลักษณะสำคัญได้ โดยน้ำหนักโมเลกุลมีอิทธิพลสูงสุด (Mean$|$SHAP$|$ = 0.1165) และสามารถจับคู่ย้อนกลับเป็นภาพโครงสร้างย่อย 2 มิติได้ 6 บิต ซึ่งมีการกระจายตัวระหว่างกลุ่ม active และ inactive ที่แตกต่างกันอย่างชัดเจน

\section{อภิปรายผลตามวัตถุประสงค์}

\begin{table}[H]
\centering
\caption{สรุปการบรรลุวัตถุประสงค์}
\label{tab:objectives}
\small
\begin{tabular}{L{5.0cm} C{1.8cm} L{6.6cm}}
\toprule
\textbf{วัตถุประสงค์} & \textbf{สถานะ} & \textbf{หลักฐานสนับสนุน}\\
\midrule
1. พัฒนาและเปรียบเทียบแบบจำลอง ML ทำนาย \pICfifty & บรรลุ & เปรียบเทียบครบ 4 แบบจำลอง บน 2 วิธีการแบ่งข้อมูล ได้ Test \Rsq\ = 0.7034\\
2. สร้างกระบวนการ Hybrid Features ผ่าน Pipeline ที่ปลอด Leakage & บรรลุ & \texttt{src/pipelines.py} ร่วมกับการทดลอง Ablation E3/E4 และการวัดผล Leakage เชิงปริมาณ E5/E5b\\
3. วิเคราะห์ SAR ด้วย SHAP และภาพโครงสร้าง 2 มิติ & บรรลุ & ตาราง Top-20 Mean$|$SHAP$|$, Beeswarm Plot และภาพโครงสร้างย่อย 6 บิต\\
\bottomrule
\end{tabular}
\end{table}

\subsection{อภิปรายวัตถุประสงค์ที่ 1}

แบบจำลองที่พัฒนาขึ้นบรรลุเกณฑ์ ``ดี'' ที่กำหนดไว้ล่วงหน้า (\Rsq\ ในช่วง 0.60--0.75) ค่า RMSE = 0.5487 log unit หมายความว่าโดยเฉลี่ยแบบจำลองทำนายค่า \ICfifty\ คลาดเคลื่อนประมาณ 3.5 เท่า ซึ่งอยู่ในระดับที่ยอมรับได้สำหรับงานคัดกรองเบื้องต้น เนื่องจากความแปรปรวนของการทดลองในห้องปฏิบัติการเองระหว่างห้องปฏิบัติการต่างกันก็อยู่ในระดับใกล้เคียงกัน

ค่า Spearman $\rho$ = 0.8415 ซึ่งสูงกว่าที่ค่า \Rsq\ บ่งชี้ เป็นข้อค้นพบที่มีนัยสำคัญเชิงปฏิบัติ เพราะในงานคัดกรองยาสิ่งที่ต้องการมากที่สุดคือความสามารถ \textbf{จัดลำดับ} สารเพื่อเลือกไปทดสอบจริง ไม่ใช่ความแม่นยำของค่าสัมบูรณ์

ประเด็นที่น่าสังเกตคือ XGBoost ทำได้ต่ำที่สุดในระยะ Cross-Validation ก่อนปรับจูน (CV \Rsq\ = 0.5803) แต่กลับดีที่สุดหลังปรับจูน (Test \Rsq\ = 0.7034) แสดงว่าค่าเริ่มต้นของ XGBoost ไม่เหมาะกับข้อมูลชุดนี้ และการปรับจูนไฮเปอร์พารามิเตอร์มีความสำคัญมากกว่าการเลือกอัลกอริทึม

\subsection{อภิปรายวัตถุประสงค์ที่ 2}

การสร้าง Pipeline ที่ป้องกัน Data Leakage บรรลุผลตามวัตถุประสงค์ในเชิงระเบียบวิธี แต่ \textbf{ผลการวัดเชิงปริมาณไม่เป็นไปตามสมมติฐานตั้งต้น}

สมมติฐานเดิมคาดว่าการปรับสเกลก่อนแบ่งข้อมูลจะทำให้ผลสูงเกินจริงอย่างเห็นได้ชัด แต่ผลการทดลอง E5 พบส่วนต่างเพียง $+0.0001$ \Rsq\ คำอธิบายคือสัดส่วนคุณลักษณะที่ต้องปรับสเกลมีเพียงร้อยละ 0.39 และขนาดชุดข้อมูลใหญ่เพียงพอ

อย่างไรก็ตาม เมื่อออกแบบการทดลองเพิ่มเติม E5b เพื่อตรวจสอบ Feature Selection Leakage ซึ่งเป็นรูปแบบที่ตรวจพบจริงในไฟล์ \texttt{xfinal.csv} ที่ได้รับมอบมา พบอคติบวก $+0.0116$ \Rsq\ ซึ่งมากกว่ารูปแบบแรกถึง 116 เท่า

\begin{keynote}[บทเรียนเชิงระเบียบวิธี]
การใช้ Pipeline ยังคงเป็นแนวปฏิบัติที่ถูกต้องและควรทำเสมอ แต่ผู้วิจัยควรตระหนักว่าความรุนแรงของ Data Leakage ขึ้นอยู่กับ \textbf{สัดส่วนของคุณลักษณะที่ได้รับผลกระทบ} และ \textbf{ขนาดของชุดข้อมูล} การรั่วไหลจากการคัดเลือกคุณลักษณะเป็นภัยคุกคามที่ใหญ่กว่าการรั่วไหลจากการปรับสเกลอย่างมีนัยสำคัญ
\end{keynote}

สำหรับผลการศึกษา Ablation พบว่า Hybrid Features ให้ผลดีที่สุดจริง แต่ส่วนเพิ่มจาก Lipinski Descriptors มีเพียง $+0.0112$ \Rsq\ ขณะที่การวิเคราะห์ SHAP กลับพบว่า Lipinski ติดอันดับ 20 แรกครบทั้ง 4 ตัว และ MW เป็นอันดับ 1 ผลสองอย่างนี้ไม่ขัดแย้งกัน แต่สะท้อนปรากฏการณ์ \textbf{คุณลักษณะสำคัญแต่ซ้ำซ้อน} เพราะ Morgan Fingerprints เข้ารหัสข้อมูลขนาดโมเลกุลไว้โดยอ้อมอยู่แล้ว

\subsection{อภิปรายวัตถุประสงค์ที่ 3}

การวิเคราะห์ SAR ให้ผลที่ตีความเชิงเคมีได้อย่างสอดคล้อง โดยเฉพาะบิต FP\_10 ที่มีอัตราส่วนความเข้มข้นระหว่างกลุ่ม active ต่อ inactive สูงถึง 9.3 เท่า และโมเลกุลตัวอย่างมีหมู่แอลดีไฮด์ซึ่งเป็นหมู่ที่ทราบกันว่าสร้างพันธะกับซิสเตอีนในช่องเร่งปฏิกิริยาของ \Mpro\ ได้

อย่างไรก็ตาม ต้องย้ำข้อจำกัดที่ประกาศไว้ตั้งแต่ต้นว่า ผลการวิเคราะห์ SHAP แสดง \textbf{ความสัมพันธ์เชิงสถิติ} ระหว่างชิ้นส่วนโมเลกุลกับค่าที่แบบจำลองทำนาย ไม่สามารถอ้างเป็น \textbf{กลไกการจับทางชีวภาพเชิงสาเหตุ} ในระดับ 3 มิติได้ ความสอดคล้องกับความรู้ทางเคมียาที่พบเป็นเพียงหลักฐานสนับสนุนความสมเหตุสมผล ไม่ใช่การพิสูจน์กลไก การยืนยันที่แท้จริงต้องอาศัย Molecular Docking, Molecular Dynamics และการทดสอบ In Vitro

กรณีบิต FP\_807 เป็นตัวอย่างที่ดีของความจำเป็นในการตรวจสอบผล SHAP ด้วยสถิติเชิงพรรณนาประกอบ เนื่องจากบิตนี้ปรากฏในโมเลกุลถึงร้อยละ 86.6 ทั้งในกลุ่ม active และ inactive ค่า SHAP ที่ติดลบจึงไม่มีความหมายเชิงเคมีที่ตีความได้ หากรายงานเฉพาะอันดับ Mean$|$SHAP$|$ โดยไม่ดูสัดส่วนการปรากฏ อาจนำไปสู่การตีความที่ผิดพลาด

\section{ข้อค้นพบสำคัญที่นอกเหนือจากวัตถุประสงค์}

\subsection{ความล้มเหลวของการทดสอบข้ามชุดข้อมูล}

ข้อค้นพบที่มีคุณค่าที่สุดของโครงงานนี้อาจเป็นผลลบ กล่าวคือ แบบจำลองที่ได้ค่า \Rsq\ = 0.7034 บนชุดทดสอบภายใน กลับได้ \Rsq\ = $-1.72$ เมื่อทดสอบกับสาร COVID-19 จากแหล่งอิสระ 90 ตัว

การวิเคราะห์ Applicability Domain อธิบายสาเหตุได้อย่างชัดเจนว่า ร้อยละ 84.4 ของสารในชุดทดสอบข้ามแหล่งมีค่าความคล้ายคลึง Tanimoto ต่อชุดฝึกต่ำกว่า 0.4 ซึ่งใกล้เคียงกับชุดสารที่ออกฤทธิ์ต่อเป้าหมายคนละชนิดโดยสิ้นเชิง

\begin{keynote}[นัยเชิงวิชาการ]
ผลนี้เป็นหลักฐานเชิงประจักษ์ที่สนับสนุนข้อวิจารณ์ที่มีต่องานวิจัย QSAR จำนวนมากที่รายงานเฉพาะค่า \Rsq\ บนชุดทดสอบภายใน โดยไม่ประเมิน Applicability Domain หรือทดสอบข้ามแหล่งข้อมูล ค่า \Rsq\ ที่สูงในการแบ่งข้อมูลแบบสุ่มไม่ได้รับประกันความสามารถทำนายในทางปฏิบัติเลย
\end{keynote}

\subsection{ผลกระทบของความไม่สมดุลของข้อมูล}

การวิเคราะห์ค่าคลาดเคลื่อนแยกตามกลุ่มฤทธิ์เผยให้เห็นว่า แบบจำลองทำนายสาร inactive สูงเกินจริงอย่างเป็นระบบเฉลี่ย 0.72 log unit ซึ่งเป็นผลจากอัตราส่วนสาร active ต่อ inactive ในชุดฝึกที่ 7.3 ต่อ 1

ข้อค้นพบนี้เชื่อมโยงโดยตรงกับความล้มเหลวในหัวข้อก่อนหน้า เพราะชุด DDH ประกอบด้วยสาร inactive ถึง 61 จาก 90 ตัว จึงตกอยู่ในสถานการณ์ที่แบบจำลองอ่อนแอที่สุดพอดี ความไม่สมดุลของข้อมูลจึงไม่ได้ส่งผลเฉพาะในงานจำแนกประเภทเท่านั้น แต่ส่งผลต่องานถดถอยด้วยเมื่อการกระจายของตัวแปรตามไม่สมมาตร

\subsection{ข้อจำกัดของ Morgan Fingerprints กับสเตอริโอเคมี}

การวิเคราะห์ Activity Cliff พบคู่โมเลกุลที่มีความคล้ายคลึงเกิน 0.85 จำนวน 635 คู่ ในจำนวนนี้มี 23 คู่ (ร้อยละ 3.62) ที่มีค่า \pICfifty\ ต่างกันเกิน 1.5 log unit

เมื่อตรวจสอบคู่ที่รุนแรงที่สุดพบว่าเป็น \textbf{สเตอริโอไอโซเมอร์} ตัวอย่างเช่น สาร CHEMBL5412571 มี \pICfifty\ = 8.52 ขณะที่ CHEMBL5197241 มี \pICfifty\ = 6.33 ทั้งสองมีสูตรโครงสร้างเชิงเส้นเหมือนกันทุกประการ ต่างกันเพียงการจัดเรียงสเตอริโอที่ศูนย์กลางไครัลสองตำแหน่ง แต่มีฤทธิ์ต่างกันถึง 2.19 log unit หรือประมาณ 155 เท่า

ค่าความคล้ายคลึง Tanimoto ระหว่างสองโมเลกุลนี้เท่ากับ \textbf{1.000 พอดี} เนื่องจาก Morgan Fingerprints ที่ตั้งค่าตามปกติไม่บันทึกข้อมูลสเตอริโอเคมี แบบจำลองจึงได้รับเวกเตอร์คุณลักษณะที่เหมือนกันทุกประการสำหรับโมเลกุลสองตัวที่มีฤทธิ์ต่างกันมหาศาล และไม่มีทางทำนายทั้งคู่ให้ถูกต้องพร้อมกันได้

ข้อค้นพบนี้มีความสำคัญมากในบริบทของสารยับยั้ง \Mpro\ เพราะสารกลุ่มนี้จำนวนมากมีลักษณะคล้ายเปปไทด์ซึ่งมีศูนย์กลางไครัลหลายตำแหน่ง และการจับกับโปรตีนขึ้นอยู่กับรูปทรงสามมิติอย่างมาก

\section{ข้อจำกัดของการศึกษา}

\begin{enumerate}[leftmargin=2em, itemsep=0.4em]
  \item \textbf{เป็นการศึกษา In Silico เท่านั้น} ไม่มีการยืนยันผลในห้องปฏิบัติการ สารที่แบบจำลองทำนายว่ามีฤทธิ์แรงยังไม่ได้รับการทดสอบจริง
  \item \textbf{SHAP แสดงความสัมพันธ์เชิงสถิติ ไม่ใช่กลไกเชิงสาเหตุ} โครงสร้างย่อยที่ระบุได้เป็นเพียงรูปแบบทางสถิติที่แบบจำลองเรียนรู้ ไม่ใช่การพิสูจน์ว่าโครงสร้างนั้นทำหน้าที่จับกับเอนไซม์จริง
  \item \textbf{ข้อมูลมาจากหลายห้องปฏิบัติการและหลายโปรโตคอล} แม้จะกรอง assay เฉพาะ \CLpro\ แล้ว แต่วิธีการวัด ความเข้มข้นของเอนไซม์ และชนิดสับสเตรตที่ต่างกันยังเป็นแหล่งความแปรปรวนในตัวแปรตาม
  \item \textbf{คุณลักษณะไม่ครอบคลุมสเตอริโอเคมี} ซึ่งเป็นข้อจำกัดเชิงหลักการที่ส่งผลต่อความสามารถทำนายสารกลุ่มเปปไทด์คล้าย
  \item \textbf{ความสามารถทำนายจำกัดอยู่ในปริภูมิเคมีของชุดฝึก} ตามที่พิสูจน์แล้วด้วยการทดสอบข้ามชุดข้อมูล
  \item \textbf{ไม่ครอบคลุม Molecular Docking และ Molecular Dynamics} ตามขอบเขตที่ประกาศไว้
\end{enumerate}

\section{ข้อเสนอแนะ}

\subsection{ข้อเสนอแนะสำหรับการนำผลไปใช้}

\begin{enumerate}[leftmargin=2em, itemsep=0.4em]
  \item \textbf{ใช้แบบจำลองเป็นเครื่องมือจัดลำดับความสำคัญ ไม่ใช่เครื่องมือวัดค่าสัมบูรณ์} เนื่องจากค่า Spearman $\rho$ สูงกว่าที่ค่า \Rsq\ บ่งชี้ การใช้แบบจำลองเพื่อเรียงลำดับสารก่อนนำไปทดสอบจริงจะให้ประโยชน์สูงสุด
  \item \textbf{ตรวจสอบ Applicability Domain ก่อนใช้ทุกครั้ง} ควรคำนวณค่า Tanimoto สูงสุดต่อชุดฝึกก่อนเชื่อผลทำนาย หากต่ำกว่า 0.4 ควรถือว่าผลทำนายไม่น่าเชื่อถือ
  \item \textbf{อ้างอิงค่า \Rsq\ = 0.5127 จาก Scaffold Split} เมื่อสื่อสารความสามารถของแบบจำลองในการทำนายสารที่มีโครงสร้างแกนกลางใหม่
\end{enumerate}

\subsection{ข้อเสนอแนะสำหรับการพัฒนาต่อยอด}

\begin{enumerate}[leftmargin=2em, itemsep=0.4em]
  \item \textbf{เปิดใช้งานสเตอริโอเคมีในการสกัด Fingerprints} โดยตั้งค่า \texttt{includeChirality=True} ซึ่งจะช่วยแก้ปัญหาสเตอริโอไอโซเมอร์ได้โดยตรง คาดว่าจะเป็นการปรับปรุงที่ให้ผลตอบแทนสูงที่สุดเมื่อเทียบกับความพยายามที่ใช้
  \item \textbf{แก้ปัญหาความไม่สมดุลของข้อมูล} ด้วยการเพิ่มน้ำหนักตัวอย่างสาร inactive ในฟังก์ชันสูญเสีย หรือรวบรวมข้อมูลสาร inactive เพิ่มเติมจากฐานข้อมูล PubChem BioAssay
  \item \textbf{ทดลองใช้คุณลักษณะเพิ่มเติม} เช่น RDKit Descriptors ชุดเต็ม, MACCS Keys หรือ Pharmacophore Fingerprints เพื่อเพิ่มมิติของสารสนเทศ
  \item \textbf{ทดลองใช้แบบจำลองที่เรียนรู้การแทนโครงสร้างเอง} เช่น Graph Neural Networks (GNN) ซึ่งรับโครงสร้างโมเลกุลเป็นกราฟโดยตรง และสามารถเข้ารหัสสเตอริโอเคมีได้
  \item \textbf{สร้างแบบจำลองประมาณความไม่แน่นอน} เช่น Conformal Prediction หรือ Ensemble Variance เพื่อให้แบบจำลองสามารถบอกได้ว่าตัวเองไม่มั่นใจเมื่อใด ซึ่งจะช่วยลดปัญหาการทำนายสารนอกโดเมน
  \item \textbf{พัฒนาเว็บแอปพลิเคชันคัดกรองตัวยา} โดยรับ SMILES เป็นข้อมูลนำเข้า แล้วแสดงค่า \pICfifty\ ที่ทำนาย พร้อมค่าความคล้ายคลึงต่อชุดฝึกเป็นตัวบ่งชี้ความน่าเชื่อถือ และภาพโครงสร้างย่อยที่มีอิทธิพลต่อการทำนาย
\end{enumerate}

\subsection{ข้อเสนอแนะเชิงระเบียบวิธีสำหรับงาน QSAR}

\begin{enumerate}[leftmargin=2em, itemsep=0.4em]
  \item \textbf{รายงานผลทั้ง Random Split และ Scaffold Split ควบคู่กันเสมอ} เพื่อให้ผู้อ่านเห็นภาพความสามารถทำนายที่แท้จริง
  \item \textbf{ประเมินและรายงาน Applicability Domain เป็นส่วนหนึ่งของผลลัพธ์มาตรฐาน} ไม่ใช่การวิเคราะห์เสริม
  \item \textbf{ตรวจสอบการซ้อนทับระหว่างชุดฝึกกับชุดทดสอบข้ามแหล่งด้วย InChIKey เสมอ} ในโครงงานนี้พบสารซ้ำ 4 ตัวซึ่งหากไม่ตัดออกจะทำให้ผลการทดสอบข้ามชุดข้อมูลไม่ถูกต้อง
  \item \textbf{วิเคราะห์ค่าคลาดเคลื่อนแยกตามกลุ่มย่อย} ไม่ใช่รายงานเฉพาะค่ารวม เพราะค่ารวมอาจซ่อนจุดอ่อนเชิงระบบไว้ ดังกรณีสาร inactive ในโครงงานนี้
  \item \textbf{ตรวจสอบผล SHAP ด้วยสถิติเชิงพรรณนาประกอบ} เช่น สัดส่วนการปรากฏของบิตในแต่ละกลุ่ม เพื่อป้องกันการตีความผิดพลาดดังกรณีบิต FP\_807
\end{enumerate}
